\documentclass[11pt]{article}
\usepackage[a4paper,margin=2.5cm]{geometry}
\usepackage{amsmath,amssymb,amsthm}
\usepackage{booktabs}
\usepackage[expansion=false]{microtype}
\usepackage{enumitem}
\usepackage[hidelinks]{hyperref}
\usepackage{xcolor}

\newcommand{\R}{\mathbb{R}}
\newcommand{\E}{\mathbb{E}}
\newcommand{\Prob}{\mathbb{P}}
\newcommand{\T}{^{\!\top}}
\newcommand{\ind}{\mathbf{1}}

\title{\textbf{A Scalable Monte-Carlo / CVaR Linear-Program Optimiser\\
for Mental-Account Portfolios with Many Securities and Derivatives}}
\author{Sami Jeddou\\[2pt]
\small Addendum to the 2012 MSc thesis (USI Lugano) and the accompanying paper\\
\small on the \emph{Beyond Mean-Variance} optimiser}
\date{\today}

\begin{document}
\maketitle

\begin{abstract}
\noindent
The thesis and its software reimplementation solve the Das--Statman mental-account
problem on an exhaustive return grid. That representation is exact for small
universes but its state space grows as $m^{N}$ in the number of securities $N$,
which confines it in practice to four or five assets and a single derivative.
This addendum sets out a \emph{scenario-based} reformulation whose cost is linear
in $N$ and independent of how returns are discretised, and which admits an
arbitrary number $K$ of derivatives priced consistently inside each scenario.
Joint returns are sampled through a copula; each mental account's
constraint---maximise expected return subject to a tail-risk floor---is then a
\emph{linear program} via the Rockafellar--Uryasev representation of Conditional
Value-at-Risk. We give the full formulation, its relationship to the grid method
and to the thesis, a convergence and validation check, and the design trade-offs.
A prototype reproduces the closed-form Normal Expected Shortfall to within $3$
basis points and optimises a five-asset, two-derivative portfolio in a single
linear solve.
\end{abstract}

\section{Motivation: the curse of dimensionality}
The Das, Markowitz, Scheid and Statman (2010) mental-account framework---extended
in the 2012 thesis with derivatives and structured products---maximises a sub-portfolio's
expected return subject to a behavioural shortfall constraint,
$\Prob(r < H) \le \alpha$. The engine evaluates this by discretising each
security's return into $m$ levels and enumerating the \emph{joint} state space.
For $N$ securities that space has on the order of $m^{N}$ cells: with $m=51$,
$N=3$ gives $\sim\!1.3\times10^{5}$ states (tractable), whereas $N=10$ gives
$51^{10}\approx10^{17}$ (infeasible). This is why exhaustive grid search is capped
at $N\le 4$ in the software and why a single derivative is the practical limit.

Institutional portfolios contain tens to hundreds of positions and multiple
hedging or yield-enhancement overlays. We therefore need a representation whose
cost grows \emph{linearly}, not exponentially, in $N$, and which scales in the
number of derivatives.

\section{A scenario-based reformulation}
Rather than enumerate a grid, draw a fixed number $S$ of joint
\emph{scenarios} from the return distribution. Each scenario $s$ is one
simultaneous realisation of all assets,
\[
   R_s = \big(R_{s,1},\dots,R_{s,N},\,R_{s,N+1},\dots,R_{s,N+K}\big)\in\R^{N+K},
\]
where the first $N$ entries are security returns and the last $K$ are derivative
returns; every scenario carries equal probability $1/S$. The data structure is an
$S\times(N+K)$ matrix, so memory and per-iteration cost are $O\!\big(S(N+K)\big)$:
\emph{linear} in the number of instruments, with $S$ chosen for statistical
accuracy and entirely decoupled from $N$. Adding a security widens each scenario
by one column; adding a derivative widens it by one more. Nothing is enumerated.

\section{Generating the scenarios}
\paragraph{Securities.}
Let $\mu\in\R^{N}$ and $\sigma\in\R^{N}$ be the horizon expected returns and
volatilities and $\Omega$ the correlation matrix, with Cholesky factor
$L$ ($\Omega = LL\T$). Drawing $\xi_s\stackrel{\text{iid}}{\sim}\mathcal N(0,I_N)$
and setting $z_s = L\xi_s$, $R_{s,i} = \mu_i + \sigma_i\,z_{s,i}$ reproduces the
thesis's multivariate-Normal assumption exactly. To capture \emph{tail dependence}
(assets falling together in stress), replace the Gaussian copula by a Student-$t$
copula: draw a multivariate-$t$ vector $t_s = z_s/\sqrt{g_s}$ with
$g_s\sim\chi^2_{\nu}/\nu$, map to uniforms $u_{s,i}=F_{t_\nu}(t_{s,i})$, and back to
the chosen marginals $z_{s,i}=\Phi^{-1}(u_{s,i})$. The dependence is then heavy-tailed
while the marginals remain controlled. This is the Sklar (1959) decomposition into
marginals and copula, standard in dependence modelling (Embrechts \emph{et al.}, 2002).

\paragraph{Derivatives (arbitrage-consistent, per scenario).}
For each derivative $k$ written on underlying $u(k)$, its entry price is the
Black--Scholes fair value $\pi_k$ at $S_0=1$ (the same convention as the engine:
strikes are fractions of the underlying, volatility is that of the underlying).
In scenario $s$ the terminal spot is $1+R_{s,u(k)}$, the payoff
$\Pi_k\!\big(1+R_{s,u(k)}\big)$ follows from the instrument's leg structure, and the
derivative's scenario return is
\[
   R_{s,N+k} \;=\; \frac{\Pi_k\!\big(1+R_{s,u(k)}\big) - \pi_k(1+p_k)}{\pi_k(1+p_k)},
\]
with $p_k$ an optional structured-product premium. Because the payoff is computed
from the \emph{same} simulated underlying that drives the security column, each
scenario keeps a coherent $(\text{underlying},\text{derivative})$ pair with no
auxiliary model. Crucially this holds for \emph{several} derivatives at once,
possibly on different underlyings: each simply contributes one more column.

\section{The optimisation problem}
For weights $w\in\R^{N+K}$ the portfolio return in scenario $s$ is the inner
product $r_s = w\T R_s$, giving $S$ weight-dependent realisations. The expected
return is the sample mean $\widehat\mu\T w$ with $\widehat\mu = \tfrac1S\sum_s R_s$.

We measure tail risk by Expected Shortfall (Conditional Value-at-Risk), the
coherent risk measure (Artzner \emph{et al.}, 1999) and the natural convex
counterpart of the thesis's loss-probability constraint. Writing the loss as
$\ell = -r$ and letting $\alpha$ be the tail probability, the
Rockafellar--Uryasev (2000, 2002) identity is
\[
   \mathrm{CVaR}_\alpha(\ell)
   \;=\; \min_{\zeta\in\R}\;\Big\{\, \zeta + \tfrac{1}{\alpha}\,
   \E\big[(\ell-\zeta)^{+}\big] \Big\},
\]
where the minimiser $\zeta^\star$ is the Value-at-Risk. On the scenario sample this
becomes $\zeta + \tfrac{1}{\alpha S}\sum_s(\ell_s-\zeta)^{+}$, which is
\emph{piecewise-linear and convex}. The tail-average return is
$\mathrm{ES}_\alpha(r) = -\mathrm{CVaR}_\alpha(\ell)$, so a behavioural floor
``the average outcome in the worst $\alpha$ tail is at least $L$'' is
$\mathrm{CVaR}_\alpha(\ell)\le -L$.

\subsection*{The linear program}
Introduce $\zeta$ (scalar) and slacks $z_s\ge 0$. The mental-account problem
\[
   \max_{w}\; \E[r] \quad\text{s.t.}\quad \mathrm{ES}_\alpha(r)\ge L,\;
   \ind\T w = 1,\; w\ge 0
\]
becomes the linear program
\[
\boxed{
\begin{aligned}
\max_{w,\zeta,z}\quad & \widehat\mu\T w \\
\text{s.t.}\quad
 & \zeta + \frac{1}{\alpha S}\sum_{s=1}^{S} z_s \;\le\; -L, \\
 & z_s \;\ge\; -\,w\T R_s - \zeta, \qquad s=1,\dots,S, \\
 & z_s \ge 0,\quad \ind\T w = 1,\quad 0 \le w \le \bar w .
\end{aligned}}
\]
The decision vector $x=(w,\zeta,z)\in\R^{\,N+K+1+S}$ has $S+1$ inequality rows; the
constraint matrix is sparse (each scenario row touches the $N{+}K$ weights, $\zeta$
and one $z_s$). Modern simplex/interior-point solvers (e.g.\ HiGHS) handle
$N+K\sim10^{2}$ and $S\sim10^{4}$ in well under a second. Convexity removes the need
for the differential-evolution and pattern-search heuristics the grid engine falls
back on beyond four assets. Optional bounds $\bar w$ (position limits), group budgets,
turnover or cardinality penalties are additional linear rows.

\section{Value-at-Risk via a CVaR proxy}
The thesis constraint $\Prob(r<H)\le\alpha$ is a \emph{chance constraint}:
non-convex and combinatorial (mixed-integer in general). Rather than solve it
exactly, we use CVaR as a conservative convex surrogate. Because
$\mathrm{CVaR}_\alpha \ge \mathrm{VaR}_\alpha$ always, imposing a CVaR floor
guarantees the VaR-style floor is met at least as safely; ES is also the coherent
measure and the one favoured by modern regulation. Where an exact
$\Prob(r<H)\le\alpha$ solution is required, the convex CVaR optimum is an excellent
warm start for a subsequent mixed-integer refinement.

\section{Multiple derivatives and mental accounts}
Two structural generalisations fall out directly:
\begin{itemize}[leftmargin=1.4em,itemsep=2pt]
\item \textbf{Several derivatives.} $K$ derivative columns are optimised jointly at
no extra modelling cost; the optimiser can, for instance, fund downside protection
on the riskiest holding while gearing upside on another. The current engine, by
contrast, prices one derivative on one underlying.
\item \textbf{Mental accounts.} Each behavioural sub-portfolio (goal) is the same
linear program over the \emph{same} scenario matrix with its own triple
$(H_j,\alpha_j,L_j)$. Solving the accounts independently and aggregating preserves
the goals-based interpretation while inheriting the scalability.
\end{itemize}

\section{Efficient frontier and warm-starting}
Sweeping the floor $L$ traces the return/tail-risk frontier. Consecutive programs
differ only in the single right-hand side $-L$, so each solve is warm-started from
the previous optimal basis and resolves almost instantly; the whole frontier costs
little more than one solve. Using common random numbers (the \emph{same} scenario
draw across the sweep) removes Monte-Carlo noise from the \emph{shape} of the
frontier.

\section{Convergence, variance reduction and validation}
Scenario estimates carry sampling error of order $S^{-1/2}$; tail quantities such as
ES need more scenarios than the distribution's body, so antithetic variates,
importance sampling of the loss tail, and common random numbers are part of the
engine rather than afterthoughts.

\paragraph{Prototype check.} A reference implementation\footnote{\texttt{mc\_cvar\_optimizer.py},
accompanying this note: scenario generator, per-scenario derivative pricing, the
CVaR linear program and a frontier sweep, in roughly 250 lines of Python/SciPy.}
on a five-asset Gaussian universe ($\alpha=5\%$, $S=4\times10^{4}$) reproduces the
closed-form Normal Expected Shortfall of an equal-weight portfolio,
\[
   \mathrm{ES}_\alpha = \mu_p - \sigma_p\,\frac{\phi\!\big(\Phi^{-1}(\alpha)\big)}{\alpha},
\]
to $3$ basis points (Table~\ref{tab:val}), returns a monotone frontier whose
realised ES sits exactly on each floor, correctly reports infeasibility when no
portfolio can reach a floor, and solves a two-derivative ($K{=}2$) instance in a
single linear program. A Student-$t$ copula deepens the equal-weight tail from
$-18.70\%$ to $-19.68\%$, as intended.

\begin{table}[h]
\centering
\begin{tabular}{lccc}
\toprule
Quantity & Monte-Carlo & Closed form & $|\Delta|$ \\
\midrule
Equal-weight $\E[r]$            & $0.0810$ & $0.0820$ & $0.0010$ \\
Equal-weight $\sigma$           & $0.1301$ & $0.1305$ & $0.0004$ \\
Equal-weight $\mathrm{ES}_{5\%}$& $-0.1870$& $-0.1873$& $0.0003$ \\
\bottomrule
\end{tabular}
\caption{Monte-Carlo estimates against closed-form Normal values
($S=4\times10^{4}$). Frontier solutions (omitted) place the realised ES on the
imposed floor to machine tolerance.}
\label{tab:val}
\end{table}

The exactness of the grid method on small universes gives a second, independent
test: the scenario optimum should converge to the grid optimum on the three- and
four-asset thesis cases, which is the recommended acceptance test before deploying
the engine at $N=100$.

\section{Relationship to the existing method and limitations}
The scenario engine \emph{complements} rather than replaces the grid. The grid is
exact and deterministic for small $N$ and remains the reference there; the scenario
engine is approximate but scales, and is the only viable route for large, multi-derivative
portfolios. The principal cautions are: (i) Monte-Carlo error, mitigated by sample
size and variance reduction; (ii) \emph{scenario quality}---the copula choice and the
estimation of $\mu,\sigma,\Omega$ drive the result, and estimation error worsens with
$N$, so shrinkage or robust estimators become part of the pipeline; and (iii) the
CVaR-for-VaR surrogacy described in \S5. None of these is a barrier to the
construction; each is a calibration responsibility.

\section{Commercial relevance}
The reformulation turns an academically faithful but small-$N$ tool into one usable
on realistic mandates: dozens to hundreds of instruments, several derivatives priced
consistently, a coherent tail-risk constraint, and a fast, warm-started frontier. The
differentiators are scenario quality, per-scenario arbitrage-consistent derivative
pricing, robustness to estimation error, and the goals-based mental-account framing
inherited from the thesis.

\begin{thebibliography}{9}
\small
\bibitem{adeh1999} Artzner, P., Delbaen, F., Eber, J.-M., Heath, D. (1999).
\emph{Coherent Measures of Risk}. Mathematical Finance, 9(3), 203--228.
\bibitem{aby2017} Alexander, G., Baptista, A., Yan, S. (2017).
\emph{Portfolio selection with mental accounts and estimation risk}.
Journal of Empirical Finance, 41, 161--186.
\bibitem{baptista2012} Baptista, A. (2012).
\emph{Portfolio selection with mental accounts and background risk}.
Journal of Banking \& Finance, 36(4), 968--980.
\bibitem{dmss2010} Das, S., Markowitz, H., Scheid, J., Statman, M. (2010).
\emph{Portfolio Optimization with Mental Accounts}.
Journal of Financial and Quantitative Analysis, 45(2), 311--334.
\bibitem{ds2009} Das, S., Statman, M. (2013).
\emph{Options and structured products in behavioral portfolios}.
Journal of Economic Dynamics and Control, 37(1), 137--153.
\bibitem{ems2002} Embrechts, P., McNeil, A., Straumann, D. (2002).
\emph{Correlation and dependence in risk management}. In: Risk Management:
Value at Risk and Beyond, Cambridge University Press, 176--223.
\bibitem{jeddou2012} Jeddou, S. (2012).
\emph{Portfolio Optimization with Mental Accounts: replication and extension with
derivatives}. MSc thesis, Universit\`a della Svizzera italiana, Lugano.
\bibitem{markowitz1952} Markowitz, H. (1952). \emph{Portfolio Selection}.
Journal of Finance, 7(1), 77--91.
\bibitem{mes2019} Momen, O., Esfahanipour, A., Seifi, A. (2019).
\emph{Collective mental accounting: an integrated behavioural portfolio selection
model for multiple mental accounts}. Quantitative Finance, 19(2), 265--275.
\bibitem{ru2000} Rockafellar, R.T., Uryasev, S. (2000).
\emph{Optimization of Conditional Value-at-Risk}. Journal of Risk, 2(3), 21--41.
\bibitem{ru2002} Rockafellar, R.T., Uryasev, S. (2002).
\emph{Conditional value-at-risk for general loss distributions}.
Journal of Banking \& Finance, 26(7), 1443--1471.
\bibitem{sklar1959} Sklar, A. (1959). \emph{Fonctions de r\'epartition \`a $n$
dimensions et leurs marges}. Publ. Inst. Statist. Univ. Paris, 8, 229--231.
\end{thebibliography}

\end{document}
